\documentclass[instructions.tex]{subfiles}
\begin{document}
 
\chapter{Il faudra expier en purgatoire les péchés qui n’auront pas été entièrement expiés en cette vie par la pénitence.}
 
Il faut, en cette matière, éviter deux écueils, l’erreur et l’illusion.
 
\primo C’est errer dans la foi, que de douter de la vérité du purgatoire : car, je demande, où va le juste qui meurt dans la grâce de Dieu, mais qui est coupable de quelques fautes légères, ou qui n’a pas entièrement satisfait à Dieu pour les fautes grièves ? Il ne peut pas être condamné à l’enfer ; parce que Dieu n’y condamne jamais le juste qui meurt dans la grâce. Il ne peut aussi être introduit dans le ciel avec ses fautes légères, parce qu’il est de foi que rien de souillé n’y entre jamais. Il faut donc qu’il y ait un lieu où cette âme achève de se purifier et de satisfaire à Dieu. Voilà ce que la foi, ce que l’Église, les saints Pères et la tradition de tous les siècles nous enseignent.
 
On doit regarder le purgatoire comme un effet de la miséricorde de Dieu, qui destine ce lieu pour être le supplément de la pénitence, et comme un effet de la justice et de la sainteté de Dieu, qui ne peut souffrir la moindre souillure, et qui se sert du feu pour épurer l’âme sainte.
 
Un chrétien éclairé craint ces feux purifians ; il ne peut penser, sans avoir le cœur pénétré de douleur, qu’il sera séparé, pour un temps, de la jouissance de son Dieu. Il regarde cependant le purgatoire avec des sentimens de reconnoissance envers Dieu, qui lui offre, dans ce lieu de tourmens, un moyen de lui satisfaire. Cette vue le porte si efficacement à aimer la bonté et la justice de Dieu, qu’il aimerait mieux souffrir, plusieurs siècles, les feux du purgatoire, que d’entrer dans le ciel, et d’y paraître devant la sainteté de Dieu avec un seul péché.
 
\secundo L’illusion la plus funeste pour ceux qui mènent une vie mondaine, c’est de s’imaginer que, s’ils ne font pas pénitence en ce monde, ils en seront quittes pour la faire en purgatoire, et qu’ensuite ils entreront dans le ciel. Détrompez-vous. L’impénitence ne conduit pas en purgatoire, mais en enfer ; puisqu’il est de foi que, sans pénitence, on sera perdu. Le purgatoire n’est pas pour commencer la pénitence, mais pour la consommer ; il n’est pas pour sanctifier l’âme, mais pour purifier l’âme sainte : il faut être pénitent, et sincèrement pénitent, pour y entrer ; et plus on est pénitent en cette vie, plus aussi en diminue-t-on la rigueur et la durée.
 
Vous vous croiriez heureux d’être en purgatoire, parce qu’il ne dure pas toujours. Oh ! que vous êtes aveugle, dit saint Augustin, de penser de la sorte ! Si les tourmens n’y sont pas éternels, ils y sont extrêmes, parce que Dieu, agissant toujours en Dieu, y punit en Dieu. C’est le même feu, dit ce saint Pere, qui purifie l’âme juste et qui brûle les réprouvés\footnote{Sub eodem igne purgatur electus et crematur damnatus. S. Aug.}.
 
On n’y souffre pas toujours, mais on y souffre quelquefois long-temps. Des fautes légères y sont punies plusieurs années ; certains péchés, qui n’ont pas été entièrement expiés ici-bas, y sont quelquefois punis des siècles entiers. Combien d’années, et peut-être des siècles, devriez-vous y souffrir, vous qui vivez dans une criminelle dissipation, et sans crainte du péché !
 
O étrange folie de s’épargner en ce monde quelques peines, et de s’exposer aux tourmens de l’autre, pour des fautes qu’il est si facile d’expier en cette vie ! Si vous étiez redevable de cent pièces d’or, et qu’on se contentât d’une seule, refuseriez-vous cette condition ? De combien de peines n’êtes-vous pas redevable à la justice de Dieu ! il serait satisfait, si vous en subissiez quelques-unes pour son amour en cette vie. Si vous le refusez, vous n’avez ni amour pour Dieu, ni charité pour vous-même.
 
Vous vous fiez aux prières qu’on fera pour vous après votre mort, aux bonnes œuvres que vous ordonnerez par votre testament ; mais il est bien plus sûr de faire vous-même vos bonnes œuvres, que de les ordonner, et d’employer saintement vos bienfaits avant la mort, que de le faire ensuite par les mains des autres.
 
Pourquoi laisser à autrui le soin de vous-même ? Ce que vous n’aurez pas fait pendant la vie, vos héritiers ne le feront pas, ou le feront mal. On est bientôt effacé du souvenir des parens et des enfans, quand on n’est plus. C’est d’ailleurs être bien imprudent de prétendre que les autres aient plus de charité pour votre âme, que vous n’en avez vous-même. Les bonnes œuvres et la pénitence sont comme nos guides pour l’éternité. Ces œuvres saintes nous conduisent bien plus sûrement quand elles précèdent que quand elles suivent. Vous prenez des mesures pour être soulagé en purgatoire ; mais prenez-les bien plus pour n’y pas tomber : il est plus sûr de préserver sa maison du feu, que de l’éteindre quand elle est embrasée. Une heure de tourmens en purgatoire est plus insupportable que tout ce qu’on peut souffrir en ce monde, et que tous les supplices des tyrans ; cependant quand il s’agit de souffrir pour expier le péché, on remet tout en l’autre vie. Quel aveuglement ! Prévenez donc par la pénitence les rigueurs du purgatoire ; mais n’oubliez pas les âmes saintes qui y sont captives ; offrez à Dieu vos prières, vos aumônes et vos bonnes œuvres, en faveur de ces âmes prédestinées, qui, après leur délivrance, seront vos intercesseurs et vos avocats auprès de Dieu.
 
\section{Résolutions}
 
\primo Je craindrai les peines du purgatoire.

\secundo Afin de les éviter, ou du moins d’en abréger la durée, je souffrirai en esprit de pénitence tout ce qui me coûtera désormais dans la pratique de mes devoirs, et dans les misères qui accompagnent cette vie.

\tertio Et pour rendre mes satisfactions efficaces, je les unirai aux souffrances de mon Sauveur. 

\prayer{O Jésus ! daignez revêtir de vos mérites mes actions et mes peines, afin qu’elles me servent à effacer peu à peu les dettes immenses que j’ai contractées auprès de votre redoutable justice, par mes nombreux péchés.}
 
\end{document}
